% 第15章第4节：加权逼近与广义Rayleigh商（广义Rayleigh商的极值、加权低秩逼近）。

\section{加权逼近与广义Rayleigh商}
\label{sec:ls-generalized-rayleigh}

前两节都用欧氏长度：单位方向按$\bm u^\top\bm u=1$规定，重建误差按各元素平方和计算。各坐标的单位或重要性不同时，长度往往按第~\ref{chap:normed-inner-product-spaces}章的加权内积度量。度量一变，有两处跟着改变：规定“单位方向”的长度变了，比较重建误差的标准也变了。

本节先处理前一处：保持二次型不变，改用加权长度规定单位方向，得到广义Rayleigh商及其极值。再处理后一处：在行、列权重可以分离的误差下研究低秩逼近，并看到它的最佳方向正是由广义Rayleigh商选出的。广义Rayleigh商还为后面的图谱松弛提供工具：卷九的谱划分把平衡约束写成按顶点度加权的长度，再用广义Rayleigh商的极值给出离散划分问题的松弛下界，可参考\BookChapterRef{chap:spectral-partition}{卷九“谱划分”章}。

\subsection{广义Rayleigh商的极值}
\label{sec:ls-generalized-eigen-reduction}
\label{sec:ls-generalized-eigen-check}

保持$\bm S=\operatorname{diag}(4,1)$不变，改用$\bm B=\operatorname{diag}(9,1)$规定长度，即要求$\bm x^\top\bm B\bm x=9x_1^2+x_2^2=1$。沿第一坐标方向，满足约束的向量是$\pm\bm e_1/3$，二次型值为$4\cdot(1/3)^2=4/9$；沿第二坐标方向，满足约束的向量仍是$\pm\bm e_2$，值为$1$。在欧氏约束下最优的$\bm e_1$方向，现在不如$\bm e_2$方向。

可见改变长度标准会改变答案。而且与示例~\ref{ex:ls-rayleigh-diagonal-two-by-two}之前的讨论一样，比较两个坐标方向还不够，必须比较全部满足加权约束的方向。与欧氏情形相同，先把约束化为一个不随缩放改变的商。
\begin{definition}[广义Rayleigh商]
\label{def:ls-generalized-rayleigh-quotient}
设$\bm S\in\mathbb R^{d\times d}$对称，$\bm B\in\mathbb R^{d\times d}$对称正定。对$\bm x\neq\bm0$，称
\begin{equation}
 R_{\bm S,\bm B}(\bm x)=\frac{\bm x^\top\bm S\bm x}{\bm x^\top\bm B\bm x}
 \label{eq:ls-generalized-rayleigh-quotient}
\end{equation}
为$\bm S$关于$\bm B$的\term{广义Rayleigh商}（generalized Rayleigh quotient）。
\end{definition}
$\bm B$正定保证分母$\bm x^\top\bm B\bm x$对一切$\bm x\neq\bm0$为正，商有定义。把$\bm x$换成$c\bm x$（$c\neq0$）时分子分母同乘$c^2$，商不变。对任意$\bm x\neq\bm0$，令$\bm u=\bm x/\sqrt{\bm x^\top\bm B\bm x}$，则$\bm u^\top\bm B\bm u=1$且$R_{\bm S,\bm B}(\bm x)=\bm u^\top\bm S\bm u$；反过来，约束$\bm x^\top\bm B\bm x=1$上的二次型值就是广义Rayleigh商。所以求广义Rayleigh商的极值，等价于在$\bm x^\top\bm B\bm x=1$上比较$\bm x^\top\bm S\bm x$。

$\bm B=\bm I$时，广义Rayleigh商恰好回到普通Rayleigh商。两者比较的是同一个二次型$\bm x^\top\bm S\bm x$，差别只在于用哪种长度规定“单位输入”：普通Rayleigh商用欧氏长度，广义Rayleigh商用加权长度$\lVert\bm x\rVert_{\bm B}$。

\paragraph{主要思想}

加权约束$\bm x^\top\bm B\bm x=1$的困难在于它不是欧氏单位球面，第~\ref{sec:ls-rayleigh-extremal}节的结论不能直接使用。思路是换一个坐标，使加权长度在新坐标中成为欧氏长度。由定理~\ref{thm:ls-gram-factor}，正定的$\bm B$可以写成$\bm B=\bm L^\top\bm L$，其中$\bm L$可逆。令$\bm z=\bm L\bm x$，则$\bm x=\bm L^{-1}\bm z$，且$\bm x\neq\bm0$当且仅当$\bm z\neq\bm0$。代入得
\[
 \bm x^\top\bm B\bm x=(\bm L\bm x)^\top(\bm L\bm x)=\bm z^\top\bm z,\qquad
 \bm x^\top\bm S\bm x=\bm z^\top\bm L^{-\top}\bm S\bm L^{-1}\bm z ,
\]
其中$\bm L^{-\top}$表示$(\bm L^{-1})^\top$。记
\[
 \bm C=\bm L^{-\top}\bm S\bm L^{-1},
\]
则$\bm C^\top=\bm L^{-\top}\bm S^\top\bm L^{-1}=\bm C$，仍是实对称矩阵，而
\[
 R_{\bm S,\bm B}(\bm x)=\frac{\bm z^\top\bm C\bm z}{\bm z^\top\bm z}=R_{\bm C}(\bm L\bm x).
\]
加权约束问题就化成了$\bm C$的普通Rayleigh商问题，可以直接应用定理~\ref{thm:matrix-analysis-rayleigh-ritz}。这一化归只用到Gram因子与换元，不需要矩阵平方根、多元微分或Lagrange乘子法。

最优方向要从新坐标换回原坐标。若$\bm C\bm z=\mu\bm z$，令$\bm v=\bm L^{-1}\bm z$，则$\bm L^{-\top}\bm S\bm v=\mu\bm L\bm v$，两边左乘$\bm L^\top$得
\[
 \bm S\bm v=\mu\bm L^\top\bm L\bm v=\mu\bm B\bm v .
\]
反过来，若$\bm S\bm v=\mu\bm B\bm v$且$\bm v\neq\bm0$，左乘$\bm L^{-\top}$并令$\bm z=\bm L\bm v$，就得到$\bm C\bm z=\mu\bm z$。所以，$\bm C$的特征对与下面方程的解一一对应。

\begin{definition}[广义特征值]
\label{def:matrix-analysis-generalized-eigenvalue}
对实对称$\bm S$与对称正定$\bm B$，称满足
\begin{equation}
\begin{gathered}
\bm S\bm v=\mu\bm B\bm v\\
\bm v\neq\bm0
\end{gathered}
\label{eq:matrix-analysis-generalized-eigenproblem}
\end{equation}
的$\mu$为广义特征值，$\bm v$为属于$\mu$的广义特征向量，求解它们的问题称为\term{广义特征问题}（generalized eigenvalue problem）。
\end{definition}

广义特征值不依赖Gram因子$\bm L$的选取。由$\bm S-\mu\bm B=\bm L^\top(\bm C-\mu\bm I)\bm L$，
\[
 \det(\bm S-\mu\bm B)=(\det\bm L)^2\det(\bm C-\mu\bm I),
\]
而$\det\bm L\neq0$，所以广义特征值恰是$\det(\bm S-\mu\bm B)=0$的根，也就是$\bm C$的特征值；不同的$\bm L$给出不同的$\bm C$，特征值却相同。

$\bm C$对称，由谱定理，这些根全是实数，按重数计共有$d$个，下面记为$\mu_1\geq\cdots\geq\mu_d$。把定理~\ref{thm:matrix-analysis-rayleigh-ritz}用于$\bm C$并换回原坐标，就得到加权约束问题的解。

\begin{theorem}[广义Rayleigh商的极值]
\label{thm:ls-generalized-rayleigh-extremal}
设$\bm S$对称、$\bm B$对称正定，广义特征值为$\mu_1\geq\cdots\geq\mu_d$。则对一切$\bm x\neq\bm0$，
\[
 \mu_d\leq R_{\bm S,\bm B}(\bm x)\leq\mu_1 .
\]
$R_{\bm S,\bm B}(\bm x)=\mu_1$当且仅当$\bm S\bm x=\mu_1\bm B\bm x$，$R_{\bm S,\bm B}(\bm x)=\mu_d$当且仅当$\bm S\bm x=\mu_d\bm B\bm x$。因此，在约束$\bm x^\top\bm B\bm x=1$上，$\bm x^\top\bm S\bm x$的最大值为$\mu_1$，最优解是满足$\bm v^\top\bm B\bm v=1$的属于$\mu_1$的广义特征向量；最小值的情形相同。
\end{theorem}
\begin{proof}
取$\bm B=\bm L^\top\bm L$与$\bm C=\bm L^{-\top}\bm S\bm L^{-1}$如上。$\bm C$的特征值就是$\mu_1\geq\cdots\geq\mu_d$。对$\bm x\neq\bm0$，令$\bm z=\bm L\bm x\neq\bm0$，由$R_{\bm S,\bm B}(\bm x)=R_{\bm C}(\bm z)$及定理~\ref{thm:matrix-analysis-rayleigh-ritz}，得$\mu_d\leq R_{\bm S,\bm B}(\bm x)\leq\mu_1$。

同一定理还说明，$R_{\bm C}(\bm z)=\mu_1$当且仅当$\bm C\bm z=\mu_1\bm z$。由上面的一一对应，这又当且仅当$\bm S\bm x=\mu_1\bm B\bm x$。$\bm C$确有属于$\mu_1$的特征向量，所以上界能够取到。最小值同理。

最后，约束$\bm x^\top\bm B\bm x=1$上的二次型值就是广义Rayleigh商，所以它的最大值为$\mu_1$，且只在满足$\bm S\bm x=\mu_1\bm B\bm x$与$\bm x^\top\bm B\bm x=1$的$\bm x$处取到。
\end{proof}

上面只用到$\bm C$两端的特征向量。对$\bm C$应用完整的谱分解，还能得到全部广义特征向量的结构：它们可以选得关于加权内积$\langle\cdot,\cdot\rangle_{\bm B}$互相正交。

\begin{theorem}[$\bm B$正交的广义特征向量]
\label{thm:ls-generalized-eigen-b-orthogonal}
设$\bm S$对称、$\bm B$对称正定，广义特征值为$\mu_1\geq\cdots\geq\mu_d$。则存在$\bm v_1,\ldots,\bm v_d$，使
\[
 \bm S\bm v_i=\mu_i\bm B\bm v_i,\qquad \bm v_i^\top\bm B\bm v_j=\delta_{i,j}\quad(i,j=1,\ldots,d).
\]
记$\bm V=[\bm v_1,\ldots,\bm v_d]$，$\bm M=\operatorname{diag}(\mu_1,\ldots,\mu_d)$，上式等价于$\bm V^\top\bm B\bm V=\bm I$与$\bm V^\top\bm S\bm V=\bm M$。
\end{theorem}
\begin{proof}
由谱定理，$\bm C=\bm Z\bm M\bm Z^\top$，其中$\bm Z=[\bm z_1,\ldots,\bm z_d]$为正交矩阵，$\bm C\bm z_i=\mu_i\bm z_i$。令$\bm V=\bm L^{-1}\bm Z$，即$\bm v_i=\bm L^{-1}\bm z_i$。由上面的一一对应，$\bm S\bm v_i=\mu_i\bm B\bm v_i$。又
\[
\begin{gathered}
 \bm V^\top\bm B\bm V=\bm Z^\top\bm L^{-\top}\bm L^\top\bm L\bm L^{-1}\bm Z=\bm Z^\top\bm Z=\bm I,\\
 \bm V^\top\bm S\bm V=\bm Z^\top\bm L^{-\top}\bm S\bm L^{-1}\bm Z=\bm Z^\top\bm C\bm Z=\bm M .
\end{gathered}
\]
逐元素比较第一式即得$\bm v_i^\top\bm B\bm v_j=\delta_{i,j}$。
\end{proof}

$\bm V$可逆，所以任意$\bm x$都能写成$\bm x=\sum_ia_i\bm v_i$。由定理，$\bm x^\top\bm B\bm x=\sum_ia_i^2$，$\bm x^\top\bm S\bm x=\sum_i\mu_ia_i^2$，广义Rayleigh商再次成为特征值的加权平均，这与普通情形完全平行，只是坐标改成了$\bm B$正交的广义特征向量。

至此两端极值都已求出。若还需要刻画中间的广义特征值，可以对$\bm C$使用第~\ref{sec:ls-spectral-value-sensitivity}节的Courant--Fischer最小最大原理（定理~\ref{thm:matrix-analysis-courant-fischer}），再按$\bm z=\bm L\bm x$换回原坐标。

\paragraph{算法实现与算例}

先回到开头的二维例子，按化归的步骤算出最优方向，并核验它们确实满足约束与广义特征方程。

\begin{example}[加权约束下的最优方向]
\label{ex:ls-generalized-rayleigh-diagonal}
取$\bm S=\operatorname{diag}(4,1)$与$\bm B=\operatorname{diag}(9,1)$，在加权约束$\bm x^\top\bm B\bm x=1$下求$\bm x^\top\bm S\bm x$的最大值、最小值与最优方向。

先化归为普通Rayleigh商。取$\bm L=\operatorname{diag}(3,1)$，则$\bm L^\top\bm L=\bm B$，
\[
 \bm C=\bm L^{-\top}\bm S\bm L^{-1}=\operatorname{diag}\Bigl(\frac49,1\Bigr).
\]
由定理~\ref{thm:matrix-analysis-rayleigh-ritz}，$R_{\bm C}$的最大值为$1$，在$\bm z=\pm\bm e_2$取到；最小值为$4/9$，在$\bm z=\pm\bm e_1$取到。换回原坐标$\bm v=\bm L^{-1}\bm z$，最大方向为$\pm\bm e_2$，最小方向为$\pm\bm e_1/3$。

再核验这两个向量满足约束与广义特征方程。由
\[
 \bm e_2^\top\bm B\bm e_2=1,\qquad \Bigl(\frac{\bm e_1}{3}\Bigr)^{\!\top}\bm B\,\frac{\bm e_1}{3}=\frac99=1,
\]
两者都满足加权单位约束。又$\bm S\bm e_2=\bm e_2=1\cdot\bm B\bm e_2$，而
\[
 \bm S\frac{\bm e_1}{3}=\frac43\bm e_1=\frac49\cdot\bm B\frac{\bm e_1}{3},
\]
所以它们分别是属于$1$与$4/9$的广义特征向量，且$\bm e_2^\top\bm B(\bm e_1/3)=0$，满足$\bm B$正交。

最后不经化归，直接比较全部可行方向。约束给出$x_2^2=1-9x_1^2$，其中$x_1^2\in[0,1/9]$，于是
\[
 \bm x^\top\bm S\bm x=4x_1^2+1-9x_1^2=1-5x_1^2\in\Bigl[\frac49,1\Bigr].
\]

最大值$1$在$x_1=0$即$\pm\bm e_2$处取到，最小值$4/9$在$x_1^2=1/9$即$\pm\bm e_1/3$处取到，与化归的结果一致。对照示例~\ref{ex:ls-rayleigh-diagonal-two-by-two}：同一个$\bm S$在欧氏约束下的最大方向是$\bm e_1$，加权约束让沿$\bm e_1$的可行向量缩短为$\bm e_1/3$，最大方向就换成了$\bm e_2$。
\end{example}

上面的例子中$\bm L$是对角矩阵，可以直接求逆。一般情形下，计算也按化归的步骤进行，但通过因子分解与线性求解完成换元，不必显式形成$\bm L^{-1}$。

第一步，用第~\ref{chap:direct-linear-solvers}章的Cholesky分解求出$\bm B$的三角因子，把$\bm L$取成上三角矩阵；若分解失败，说明$\bm B$不是正定矩阵，正定化归不能继续使用。第二步，分两次三角求解形成$\bm C$：先解$\bm L^\top\bm W=\bm S$，得到$\bm W=\bm L^{-\top}\bm S$；再解$\bm L^\top\bm C^\top=\bm W^\top$，得到$\bm C=\bm W\bm L^{-1}$。

第三步，对对称矩阵$\bm C$作正交谱分解，按所求的是最大值还是最小值取单位特征向量$\bm z$。第四步，解三角方程组$\bm L\bm v=\bm z$换回原坐标；由$\bm v^\top\bm B\bm v=\bm z^\top\bm z=1$，所得$\bm v$自动满足加权单位约束。最后核验广义特征残差$\bm S\bm v-\mu\bm B\bm v$与加权长度$\bm v^\top\bm B\bm v$。整个过程保持对称结构，只需对称矩阵的特征值算法。

\begin{example}[按三角求解的步骤计算]
\label{ex:matrix-approximation-generalized-cholesky}
取$\bm S=\operatorname{diag}(2,3)$、$\bm B=\operatorname{diag}(4,1)$。$\bm B$的Cholesky因子为$\bm L=\operatorname{diag}(2,1)$。解$\bm L^\top\bm W=\bm S$得$\bm W=\operatorname{diag}(1,3)$，再由$\bm C=\bm W\bm L^{-1}$得$\bm C=\operatorname{diag}(1/2,3)$。

$\bm C$的最大特征值$3$对应$\bm z=\bm e_2$，解$\bm L\bm v=\bm z$得$\bm v=(0,1)^\top$；最小特征值$1/2$对应$\bm z=\bm e_1$，得$\bm v=(1/2,0)^\top$。核验：$\bm S(0,1)^\top=(0,3)^\top=3\bm B(0,1)^\top$，$\bm S(1/2,0)^\top=(1,0)^\top=\frac12\cdot(2,0)^\top=\frac12\bm B(1/2,0)^\top$，且$(1/2)^2\cdot4=1$。所以加权约束下的最大值为$3$、最小值为$1/2$。

与欧氏约束比较：$\bm S$的普通Rayleigh商最大为$3$、最小为$2$，最优直线同样是两条坐标轴。这里加权没有改变最优直线，却把最小值由$2$变成$1/2$，因为沿$\bm e_1$满足加权约束的向量缩短为$\bm e_1/2$。与示例~\ref{ex:ls-generalized-rayleigh-diagonal}合起来看，长度标准既可能改变最优方向，也可能只改变最优值。
\end{example}

\paragraph{算法分析}

形成$\bm C$只用到Cholesky分解与以三角矩阵为系数的求解，代价都是$O(d^3)$次运算；之后的主要代价在对称特征值计算，所用算法本书不作分析。只需要最大或最小的广义特征对时，也必须先完成同样的换元，因为$\bm S$本身的特征向量一般不是广义特征向量。

另一种看似直接的做法是把$\bm S\bm v=\mu\bm B\bm v$左乘$\bm B^{-1}$，改求普通特征问题$\bm B^{-1}\bm S\bm v=\mu\bm v$。由$\bm B^{-1}\bm S=\bm L^{-1}\bm L^{-\top}\bm S=\bm L^{-1}\bm C\bm L$，它与$\bm C$相似，特征值相同；但$\bm B^{-1}\bm S$一般不对称。

例如$\bm S=\left[\begin{smallmatrix}2&1\\1&2\end{smallmatrix}\right]$、$\bm B=\operatorname{diag}(2,1)$时，$\bm B^{-1}\bm S=\left[\begin{smallmatrix}1&1/2\\1&2\end{smallmatrix}\right]$。它的特征多项式为$z^2-3z+\tfrac32$，特征值为$\mu_{1,2}=(3\pm\sqrt3)/2$；由第一行$(1-\mu)x_1+\tfrac12x_2=0$得特征向量$\bm v_1=(1,1+\sqrt3)^\top$、$\bm v_2=(1,1-\sqrt3)^\top$。这两个向量就是广义特征向量，满足$\bm v_1^\top\bm B\bm v_2=2+(1-3)=0$，关于$\bm B$正交；但$\bm v_1^\top\bm v_2=1+(1-3)=-1\neq0$，在欧氏内积下并不正交。所以不能对$\bm B^{-1}\bm S$套用谱定理与欧氏正交特征基的结论。前面两个二维例子里$\bm B^{-1}\bm S$恰好是对角矩阵，只是特殊情形。

本小节的化归依赖$\bm B$正定。若$\bm B$不是正定矩阵，$\bm x^\top\bm B\bm x$不再定义长度：$\bm B$半正定而奇异时，存在$\bm x\neq\bm0$使分母为零；$\bm B$不定时，$\bm x^\top\bm B\bm x$可以为负。这时也不存在可逆的$\bm L$使$\bm B=\bm L^\top\bm L$：若有这样的$\bm L$，则$\bm x\neq\bm0$时$\bm L\bm x\neq\bm0$，$\bm x^\top\bm B\bm x=\lVert\bm L\bm x\rVert_2^2>0$，$\bm B$就是正定的。

于是化归无从进行，约束问题甚至可能没有极值。例如$\bm S=\bm I$、$\bm B=\operatorname{diag}(1,0)$时，约束$x_1^2=1$不限制$x_2$，而$\bm x^\top\bm S\bm x=1+x_2^2$没有上界。这类更一般的广义特征问题不在此展开。

\subsection{加权低秩逼近}
\label{sec:matrix-approximation-weighted-low-rank}

\paragraph{主要思想}

现在改变重建误差的度量。设$\bm A\in\mathbb R^{m\times d}$的各行表示不同的量，重要性或单位不同；各列表示不同的样本，可信程度也不同。给定对称正定的行权重$\bm\Omega_{\mathrm r}\in\mathbb R^{m\times m}$与列权重$\bm\Omega_{\mathrm c}\in\mathbb R^{d\times d}$，用
\[
 E(\bm X)=\operatorname{tr}\bigl((\bm A-\bm X)^\top\bm\Omega_{\mathrm r}(\bm A-\bm X)\bm\Omega_{\mathrm c}\bigr)
\]
度量$\bm X$的误差，并在$\operatorname{rank}\bm X\leq k$下求其最小值，称为\term{加权低秩逼近}（weighted low-rank approximation）。$\bm\Omega_{\mathrm r}=\bm I$、$\bm\Omega_{\mathrm c}=\bm I$时它就是上一节的Frobenius平方误差。

两个权重都是对角矩阵时，记$\bm\Omega_{\mathrm r}=\operatorname{diag}(\omega_1,\ldots,\omega_m)$、$\bm\Omega_{\mathrm c}=\operatorname{diag}(\tau_1,\ldots,\tau_d)$，展开迹得
\[
 E(\bm X)=\sum_{i=1}^m\sum_{j=1}^d\omega_i\tau_j(a_{i,j}-x_{i,j})^2 ,
\]
即第$(i,j)$个元素的平方偏差乘以权重$\omega_i\tau_j$。所以这一类误差包含逐元素加权，但要求权重能分离成行因子与列因子的乘积。

为了化归，取Gram因子。由定理~\ref{thm:ls-gram-factor}，存在可逆的$\bm K$与$\bm N$，使$\bm\Omega_{\mathrm r}=\bm K^\top\bm K$、$\bm\Omega_{\mathrm c}=\bm N\bm N^\top$（对$\bm\Omega_{\mathrm c}$取$\bm N^\top$为它的Gram因子）。记$\bm M=\bm A-\bm X$，由迹的循环性，
\[
 \lVert\bm K\bm M\bm N\rVert_{\mathrm F}^2
 =\operatorname{tr}(\bm N^\top\bm M^\top\bm K^\top\bm K\bm M\bm N)
 =\operatorname{tr}(\bm M^\top\bm\Omega_{\mathrm r}\bm M\bm N\bm N^\top)
 =E(\bm X).
\]
于是$E(\bm X)=\lVert\bm K(\bm A-\bm X)\bm N\rVert_{\mathrm F}^2$：加权误差是变换后矩阵的普通Frobenius误差。

令$\widetilde{\bm A}=\bm K\bm A\bm N$、$\bm Z=\bm K\bm X\bm N$，则$E(\bm X)=\lVert\widetilde{\bm A}-\bm Z\rVert_{\mathrm F}^2$。秩约束在这个变换下不变：乘积的列空间含于左因子的列空间，所以$\operatorname{rank}\bm Z\leq\operatorname{rank}\bm X$；又$\bm X=\bm K^{-1}\bm Z\bm N^{-1}$，同理$\operatorname{rank}\bm X\leq\operatorname{rank}\bm Z$，故二者相等。

于是$\bm X\mapsto\bm K\bm X\bm N$把秩不超过$k$的矩阵一一地变为秩不超过$k$的矩阵，原问题等价于对$\widetilde{\bm A}$作普通的低秩逼近。由定理~\ref{thm:ls-best-low-rank}与定理~\ref{thm:ls-best-low-rank-uniqueness}立即得到下面的结论。

\begin{theorem}[可分离权重下的最佳低秩逼近]
\label{thm:matrix-approximation-weighted-low-rank}
设$\bm\Omega_{\mathrm r}=\bm K^\top\bm K$、$\bm\Omega_{\mathrm c}=\bm N\bm N^\top$，$\bm K$、$\bm N$可逆，$\widetilde{\bm A}=\bm K\bm A\bm N$的奇异值为$\tilde\sigma_1\geq\cdots\geq\tilde\sigma_r>0$，其中$r=\operatorname{rank}\widetilde{\bm A}=\operatorname{rank}\bm A$。对$0\leq k<r$，
\[
 \min_{\operatorname{rank}\bm X\leq k}E(\bm X)=\sum_{i=k+1}^r\tilde\sigma_i^2 ,
\]
最小值由$\bm X_k=\bm K^{-1}\widetilde{\bm A}_k\bm N^{-1}$达到，其中$\widetilde{\bm A}_k$是$\widetilde{\bm A}$的截断SVD。若$1\leq k<r$且$\tilde\sigma_k>\tilde\sigma_{k+1}$，最优解唯一。
\end{theorem}
\begin{proof}
由上面的等式，$E(\bm X)=\lVert\widetilde{\bm A}-\bm K\bm X\bm N\rVert_{\mathrm F}^2$，而$\bm X\mapsto\bm K\bm X\bm N$是秩不超过$k$的矩阵集合到自身的一一映射。所以$E$在原集合上的最小值等于$\lVert\widetilde{\bm A}-\bm Z\rVert_{\mathrm F}^2$在$\operatorname{rank}\bm Z\leq k$上的最小值，由定理~\ref{thm:ls-best-low-rank}，它等于$\sum_{i>k}\tilde\sigma_i^2$，在$\bm Z=\widetilde{\bm A}_k$即$\bm X=\bm X_k$处取到。有间隙时，由定理~\ref{thm:ls-best-low-rank-uniqueness}，最优的$\bm Z$唯一，一一映射保证最优的$\bm X$也唯一。
\end{proof}

最优解的结构可以用上一小节的广义Rayleigh商看清。只设行权重，取$\bm\Omega_{\mathrm c}=\bm I$、$\bm N=\bm I$，并取$k=1$。变换后的秩一候选写成$\bm Z=\bm z\bm c^\top$，$\lVert\bm z\rVert_2=1$；令$\bm w=\bm K^{-1}\bm z$，则原坐标中的候选是$\bm X=\bm w\bm c^\top$，单位约束变为$\bm w^\top\bm\Omega_{\mathrm r}\bm w=\bm z^\top\bm z=1$。

按第~\ref{sec:ls-rank-one-rayleigh}节，变换后问题的保留量是
\[
 \bm z^\top\widetilde{\bm A}\widetilde{\bm A}^\top\bm z
 =\bm w^\top\bm K^\top\bm K\bm A\bm A^\top\bm K^\top\bm K\bm w
 =\bm w^\top\bm\Omega_{\mathrm r}\bm A\bm A^\top\bm\Omega_{\mathrm r}\bm w .
\]
所以最佳方向$\bm w$使广义Rayleigh商$R_{\bm\Omega_{\mathrm r}\bm A\bm A^\top\bm\Omega_{\mathrm r},\bm\Omega_{\mathrm r}}$取最大值，是广义特征问题$\bm\Omega_{\mathrm r}\bm A\bm A^\top\bm\Omega_{\mathrm r}\bm w=\mu\bm\Omega_{\mathrm r}\bm w$的最大特征向量。按本节的化归，这里的$\bm C=\bm K^{-\top}(\bm\Omega_{\mathrm r}\bm A\bm A^\top\bm\Omega_{\mathrm r})\bm K^{-1}=\widetilde{\bm A}\widetilde{\bm A}^\top$，最大广义特征值就是$\tilde\sigma_1^2$。

对应的最优系数是$\bm c=\widetilde{\bm A}^\top\bm z=\bm A^\top\bm\Omega_{\mathrm r}\bm w$，最优近似为$\bm X_1=\bm w\bm w^\top\bm\Omega_{\mathrm r}\bm A$。矩阵$\bm w\bm w^\top\bm\Omega_{\mathrm r}$满足$(\bm w\bm w^\top\bm\Omega_{\mathrm r})^2=\bm w\bm w^\top\bm\Omega_{\mathrm r}$，且$\bm\Omega_{\mathrm r}\bm w\bm w^\top\bm\Omega_{\mathrm r}$对称，它是关于加权内积$\langle\cdot,\cdot\rangle_{\bm\Omega_{\mathrm r}}$到$\operatorname{span}\{\bm w\}$的正交投影。所以加权秩一逼近与第~\ref{sec:ls-rank-one-rayleigh}节完全平行，只是长度与投影都换成了按$\bm\Omega_{\mathrm r}$度量。

权重改变了用于排序的奇异方向：决定保留哪些方向的是$\widetilde{\bm A}$的奇异值，而不是$\bm A$的。因此直接截断原矩阵$\bm A$，一般不再是加权意义下的最优解，下面的例子给出具体比较。

这一化归要求误差能写成左右乘法的形式，任意的逐元素权重未必满足。若对一切$\bm X$都有$E(\bm X)=\sum_{i,j}w_{i,j}(a_{i,j}-x_{i,j})^2$，取$\bm A-\bm X=\bm e_i\bm e_j^\top$，左端为$(\bm\Omega_{\mathrm r})_{i,i}(\bm\Omega_{\mathrm c})_{j,j}$，右端为$w_{i,j}$，于是权重矩阵必是行因子与列因子的乘积，$w_{i,j}w_{k,l}=w_{i,l}w_{k,j}$。例如$\left[\begin{smallmatrix}1&2\\2&1\end{smallmatrix}\right]$不满足$1\cdot1=2\cdot2$，不能这样化归；缺失元素对应权重为零，也不满足正定性要求。

这些情形下误差不再按任何正交投影分解，没有截断SVD式的闭式解，一般只能用迭代优化求解，并且要分别讨论所得解是否为全局最优。这类问题的通用求解方法由卷五展开，可参考\BookChapterRef{chap:optimization}{卷五“优化问题与解”章}。

\paragraph{算法实现与算例}

计算按化归的顺序进行。第一步，对$\bm\Omega_{\mathrm r}$与$\bm\Omega_{\mathrm c}$作Cholesky分解得到三角因子$\bm K$、$\bm N$；权重是对角矩阵时直接取各对角元的平方根。第二步，形成$\widetilde{\bm A}=\bm K\bm A\bm N$，求它的SVD并截断，得到$\widetilde{\bm A}_k$。

第三步，还原$\bm X_k=\bm K^{-1}\widetilde{\bm A}_k\bm N^{-1}$：先解$\bm Y\bm N=\widetilde{\bm A}_k$，即$\bm N^\top\bm Y^\top=\widetilde{\bm A}_k^\top$，再解$\bm K\bm X_k=\bm Y$，用两次三角求解代替显式求逆。若只需要低秩因子，可以只对$\widetilde{\bm U}_k$与$\widetilde{\bm V}_k^\top$分别作这两次求解，得到$\bm X_k=(\bm K^{-1}\widetilde{\bm U}_k)\widetilde{\bm\Sigma}_k(\widetilde{\bm V}_k^\top\bm N^{-1})$，秩不超过$k$由因子形式直接保证。最后在原坐标下核验秩与加权残差$E(\bm X_k)=\sum_{i>k}\tilde\sigma_i^2$。

\begin{example}[行加权改变保留的方向]
\label{ex:matrix-approximation-weighted-low-rank}
取$\bm A=\operatorname{diag}(3,2)$，行权重$\bm\Omega_{\mathrm r}=\operatorname{diag}(1,4)$，列权重$\bm\Omega_{\mathrm c}=\bm I$，$k=1$。取$\bm K=\operatorname{diag}(1,2)$、$\bm N=\bm I$，变换后为$\widetilde{\bm A}=\operatorname{diag}(3,4)$，奇异值为$4$（方向$\bm e_2$）与$3$（方向$\bm e_1$）。

保留第二方向得$\widetilde{\bm A}_1=\operatorname{diag}(0,4)$，还原为$\bm X_1=\bm K^{-1}\widetilde{\bm A}_1=\operatorname{diag}(0,2)$。它的加权误差是$\lVert\bm K(\bm A-\bm X_1)\rVert_{\mathrm F}=\lVert\operatorname{diag}(3,0)\rVert_{\mathrm F}=3=\tilde\sigma_2$。直接截断$\bm A$得到$\operatorname{diag}(3,0)$，加权误差为$\lVert\bm K\operatorname{diag}(0,2)\rVert_{\mathrm F}=\lVert\operatorname{diag}(0,4)\rVert_{\mathrm F}=4$，更大。

用广义Rayleigh商核对：$\bm\Omega_{\mathrm r}\bm A\bm A^\top\bm\Omega_{\mathrm r}=\operatorname{diag}(9,64)$，关于$\bm\Omega_{\mathrm r}=\operatorname{diag}(1,4)$的广义特征值为$9$与$16$。最大者$16=\tilde\sigma_1^2$对应$\bm w=\bm e_2/2$（满足$\bm w^\top\bm\Omega_{\mathrm r}\bm w=1$），系数$\bm c=\bm A^\top\bm\Omega_{\mathrm r}\bm w=(0,4)^\top$，重建$\bm w\bm c^\top=\operatorname{diag}(0,2)=\bm X_1$。第二行的权重放大了它的偏差，所以虽然$\bm A$的第一列更长，最优近似却保留了第二个坐标方向。
\end{example}

\paragraph{算法分析}

除一次$m\times d$矩阵的SVD外，其余步骤是两次Cholesky分解、两次矩阵乘法与若干三角求解，代价为$O(m^3+d^3+md(m+d))$次运算。权重必须正定；某行或某列权重为零时，相应的$\bm K$或$\bm N$不可逆，还原步骤无法进行。

还原步骤还会放大计算误差。若$\widetilde{\bm A}_k$带有误差$\bm\Delta$，还原后的误差为$\bm K^{-1}\bm\Delta\bm N^{-1}$；逐列、逐行应用谱范数的定义，可得$\lVert\bm K^{-1}\bm\Delta\bm N^{-1}\rVert_{\mathrm F}\leq\lVert\bm K^{-1}\rVert_2\lVert\bm\Delta\rVert_{\mathrm F}\lVert\bm N^{-1}\rVert_2$。权重的最大、最小特征值相差悬殊时，这个放大倍数可能很大，原坐标下的结果需要按第~\ref{sec:matrix-approximation-verification}节的方法核验。
